\section{Domain-Driven Design}

Domain-Driven Design (DDD) tổ chức phần mềm quanh ngôn ngữ và quy tắc nghiệp
vụ thay vì quanh bảng dữ liệu hoặc framework \cite{evans2003ddd}. Cách tiếp cận
này phù hợp với nền tảng có nhiều miền như Fleet, Model Registry, Deployment,
Incident, Knowledge, Copilot và Ops.

\subsection{Entity}

Entity có định danh ổn định và vòng đời thay đổi theo thời gian. Ví dụ gồm Node,
Model, Deployment, Command và Incident. Hai entity có cùng thuộc tính nhưng
khác định danh vẫn là hai đối tượng khác nhau.

\subsection{Value Object}

Value Object được xác định bởi giá trị và thường bất biến, chẳng hạn ModelVersion,
Checksum, DeviceCapability, MetricWindow hoặc DeploymentPolicy. Value Object
giúp đóng gói validation và tránh truyền các chuỗi rời rạc không có ngữ nghĩa.

\subsection{Aggregate}

Aggregate là biên nhất quán cho một nhóm entity và value object. Aggregate root
kiểm soát mọi thay đổi bên trong. Deployment có thể là aggregate quản lý target,
batch, trạng thái và policy; các thay đổi phải đi qua hành vi của Deployment thay
vì cập nhật bảng trực tiếp.

\subsection{Repository}

Repository cung cấp giao diện tải và lưu aggregate theo ngôn ngữ domain. Domain
không phụ thuộc chi tiết PostgreSQL, object storage hoặc message broker. Điều này
giúp kiểm thử quy tắc nghiệp vụ độc lập với hạ tầng.

\subsection{Domain Service}

Domain Service chứa hành vi nghiệp vụ không thuộc tự nhiên về một entity, ví dụ
chọn nhóm canary dựa trên capability hoặc đánh giá chính sách triển khai. Service
không nên trở thành nơi chứa toàn bộ logic; invariant vẫn được giữ trong aggregate
khi có thể.

\subsection{Bounded Context}

Bounded Context xác định biên mô hình và ngôn ngữ. Fleet quản lý node và
capability; Model Registry quản lý model/version/artifact; Deployment quản lý
rollout; Incident quản lý sự cố; Knowledge quản lý tài liệu; Copilot điều phối hội
thoại và tool call; Ops cung cấp quan sát vận hành. Các context trao đổi qua API
hoặc event với hợp đồng rõ ràng.